\newpage
\section{量子比特路由}

\subsection{问题定义}

量子比特路由（Quantum Circuit Routing）是将逻辑量子电路映射到物理量子硬件的过程。由于 NISQ 设备中只有相邻的物理量子比特可以执行双量子比特门（如 CNOT），当逻辑门操作的两个量子比特在物理上不相邻时，必须插入 SWAP 门来交换量子比特的位置。每次 SWAP 操作引入额外噪声并增加电路深度，因此目标是在满足物理约束的前提下，最小化 SWAP 数量或最大化最终保真度。

更形式化地，给定：
\begin{itemize}
  \item 逻辑电路 $C$ 包含 $n$ 个逻辑量子比特和一系列门操作
  \item 物理拓扑 $G = (V, E)$ 包含 $m \geq n$ 个物理量子比特，$E$ 为相邻边集
  \item 初始映射 $\pi: [n] \to [m]$ 将逻辑量子比特映射到物理位置
\end{itemize}

每一步，智能体从 $E$ 中选择一条边 $(p,q)$ 执行 SWAP 操作，这一操作交换物理量子比特 $p$ 和 $q$ 上的内容。在 SWAP 后，所有前置门已执行完的双量子比特门自动执行。重复此过程，直到电路中的所有门都执行完毕或达到最大步数限制。

\subsection{总体架构}

系统采用强化学习结合图神经网络的方法解决路由问题：GNN 编码器将电路结构和硬件拓扑编码为节点嵌入，PPO 智能体根据这些嵌入为每条候选 SWAP 边打分，选取得分最高的动作执行。

\begin{figure}[h]
\centering
\begin{verbatim}
  NoiseSimulator (sim/sim.py)
       │
       └──> HardwareFeatures (features.py) ──> CircuitDAG (circuit_dag.py)
              │                                        │
              └──> SubGNN (encoder.py)                 │
                     │                                 │
                     ├──> RoutingEnv (env.py)          │
                     │        │                        │
                     └──> PPOAgent (agent.py) ─────────┘
                              │
                              ├──> train_agent.py (training)
                              └──> eval_policy.py (inference)
\end{verbatim}
\caption{系统架构图}
\label{fig:arch}
\end{figure}

\subsubsection{CircuitDAG（电路有向无环图）}

\texttt{routing/graph/circuit\_dag.py} 将 QuantumCircuit 解析为 DAG 结构，每个节点对应一个门操作，边表示依赖关系。每个门记录其前驱门索引，支持拓扑排序和 front\_layer 提取（所有依赖条件已满足且未执行的双量子比特门）。特征编码输出 32 维节点特征和 16 维边特征。

\subsubsection{HardwareFeatures（硬件特征）}

\texttt{routing/graph/features.py} 从噪声配置构建硬件特征：包含物理距离矩阵、噪声参数（T1/T2、门错误率、串扰强度等），并提供 \texttt{RoutingGraphData} 数据结构作为 GNN 的 PyTorch Geometric Data 对象。

\subsubsection{SubGNN（图神经网络编码器）}

\texttt{routing/gnn/encoder.py} 实现基于 GATConv 的双向二分图编码器：
\begin{itemize}
  \item \textbf{输入}：32 维节点特征（逻辑/物理量子比特）、16 维边特征（门操作/物理连接）
  \item \textbf{架构}：3 层多头 GAT (hidden=64, heads=4) + LayerNorm + ReLU
  \item \textbf{输出}：\texttt{node\_embeddings()} 返回每个物理量子比特的 48 维嵌入
  \item \textbf{全局池化}：mean+max 池化得到 96 维全局嵌入（仅用于历史对比）
\end{itemize}

\subsubsection{RoutingEnv（强化学习环境）}

\texttt{routing/rl/env.py} 实现 Gymnasium 接口的环境，核心设计如下：

\textbf{观测空间}：每条耦合边独立构建局部特征向量：
\[
\text{edge\_feat}_i = [h_p, h_q, h_p-h_q, \text{sabre\_feats}] \in \mathbb{R}^{149}
\]
其中 $h_p, h_q \in \mathbb{R}^{48}$ 为 GNN 节点嵌入，$\text{sabre\_feats} \in \mathbb{R}^5$ 为 SABRE 启发式特征。完整观测为所有边的特征拼接，加上映射向量和进度标量。

\textbf{动作空间}：\texttt{Discrete(num\_edges)}，每个动作选择一条物理耦合边执行 SWAP。

\textbf{奖励函数}（Phase 1）：
\[
r_t = r_{\text{swap}} + \sum r_{\text{execute}} + \sum r_{\text{propagate}} + r_{\text{dist}}
\]
\begin{itemize}
  \item $r_{\text{swap}} = -0.3$：每次 SWAP 的固定代价
  \item $r_{\text{execute}}$：门执行奖励（CX=2.0，1Q=0.3--0.5）
  \item $r_{\text{propagate}}$：X/Z 错误传播惩罚
  \item $r_{\text{dist}} = -\eta \cdot (d_{\text{after}} - d_{\text{before}}) / \max(d_{\text{before}}, 1)$：距离缩短奖励（$\eta = 1.0$）
\end{itemize}

Phase 2 增加终端保真度奖励：$R_{\text{terminal}} = \lambda_{\text{fid}} \cdot F$，其中 $F$ 通过 Aer 密度矩阵仿真计算。

\textbf{死锁掩码}：检测最近 2 步内同一边的重复 SWAP 震荡，将该边屏蔽以防止无限循环。

\textbf{克隆方法}：\texttt{clone()} 返回环境的深拷贝，用于 beam search 推理时的状态模拟。

\subsubsection{PPOAgent（策略智能体）}

\texttt{routing/rl/agent.py} 实现基于 PPO 的智能体：

\textbf{EdgeActorCritic 网络}：
\[
\text{score}_i = \text{MLP}_{\text{actor}}(\text{edge\_feat}_i) \in \mathbb{R}
\]
\[
\pi(a|s) = \text{softmax}([\text{score}_0, \ldots, \text{score}_{K-1}])
\]
Critic 使用注意力池化聚合边特征，降低输入维度至 155，提高训练稳定性。

\textbf{超参数}：$\text{lr}=3\times10^{-4}$, $\gamma=0.99$, $\lambda=0.95$, $\text{clip}=0.2$, $\text{entropy}=0.01$

\textbf{联合训练}：GNN 编码器参数纳入 PPO 优化器，实现 RL 驱动的端到端特征学习。

\subsection{训练流水线}

系统采用两阶段课程训练策略：

\subsubsection{Phase 1：纯路由训练}

目标：学习高效路由策略，最小化 SWAP 数。
\begin{itemize}
  \item 时间步数：100,000
  \item 奖励模式：\texttt{routing}（无终端保真度信号）
  \item SABRE 距离特征 + 距离即时奖励 + 死锁掩码
  \item 数据集课程：Phase1 $\to$ Phase2 $\to$ Phase3（深度 2--13）
\end{itemize}

训练日志显示 entropy 从 1.38 降至 0.79（约 45K 步收敛），value loss 降至 0.1，SWAP 均值稳定在 4.6。

\subsubsection{Phase 2：噪声感知微调}

目标：在保持路由能力的同时优化保真度。
\begin{itemize}
  \item 加载 Phase 1 检查点
  \item 时间步数：100,000
  \item 奖励模式：\texttt{noise\_aware}（Aer 仿真保真度终端奖励）
  \item $\lambda_{\text{fid}}$ 从 0 线性退火至 5.0（前 50\% 步）
  \item 多拓扑混合：cross\_5q, ring\_5q, ibmq\_5\_line
\end{itemize}

全程 trunc=0\%，entropy 稳定在 0.79，value loss 0.06--0.29，fidelity 约 0.78。

\subsection{Beam Search 推理}

为了弥补 PPO 前馈推理无法搜索的缺陷，实现了 1 步 Beam Search：
\begin{enumerate}
  \item 策略网络输出所有边的 logit 分数
  \item 选取得分最高的 K 个动作（默认 K=3）
  \item 对每个候选动作：克隆环境、执行虚拟 step、用 Critic 评估 $V(s')$
  \item 选取得分最高的动作执行真实 step
\end{enumerate}

\subsection{测试情况}

当前共有 3 个测试文件，18 个测试用例全部通过：

\begin{itemize}
  \item \textbf{test\_circuit\_dag.py}（5 个测试）：DAG 解析、图特征张量形状、串扰特征、PyG Data 转换
  \item \textbf{test\_env.py}（12 个测试）：观测空间、动作执行、奖励计算、噪声感知模式、死锁掩码、错误追踪
  \item \textbf{test\_eval\_phase1.py}（1 个慢测试）：端到端模型评估，验证 PPO 在三个拓扑上 100\% 完成，SWAP 数低于 Greedy
\end{itemize}

\subsection{当前最佳结果}

\subsubsection{纯路由模式（routing mode, deterministic）}

\begin{center}
\begin{tabular}{|c|c|c|c|c|}
\hline
拓扑 & PPO argmax & PPO beam3 & SABRE & Greedy \\
\hline
cross\_5q & 2.0 & \textbf{1.3} & \textbf{1.2} & 1.8 \\
ring\_5q & 1.6 & \textbf{1.2} & \textbf{1.0} & 1.1 \\
ibmq\_5\_line & 2.8 & \textbf{2.3} & \textbf{2.0} & 3.1 \\
\hline
\end{tabular}
\end{center}

\subsubsection{噪声感知模式（noise\_aware mode, deterministic）}

\begin{center}
\begin{tabular}{|c|c|c|c|c|}
\hline
拓扑 & PPO SWAPs & SABRE SWAPs & PPO 保真度 & SABRE 保真度 \\
\hline
cross\_5q & 3.2 & 3.0 & \textbf{0.9453} & 0.9327 \\
ring\_5q & 4.5 & 3.5 & \textbf{0.9419} & 0.9371 \\
ibmq\_5\_line & 6.7 & 5.5 & \textbf{0.9557} & 0.9228 \\
\hline
\end{tabular}
\end{center}

关键发现：PPO 在 line/cross 上的保真度已超过 SABRE（分别高出 3.29pp 和 1.26pp），Beam Search 将路由 SWAP 差距闭合了 63--88\%。

\subsection{下一步研究方向}

\subsubsection{多步 Beam Search 与 MCTS}

当前 1 步 lookahead 在线性拓扑上闭合率最低（63\%），因为该拓扑需要多步规划。扩展到 beam depth=2 或 3，并在叶节点进行完整的保真度评估，有望进一步缩小差距。中长期可以引入 MCTS 作为推理时的搜索策略（类似 AlphaZero 范式）。

\subsubsection{更大规模拓扑与电路}

当前系统仅支持 5 量子比特拓扑，实际 NISQ 设备可达 27--127 量子比特。主要挑战：
\begin{itemize}
  \item 动作空间从 4--5 条边增长到 30--100+ 条边，策略网络需要扩展到处理可变大边集
  \item GNN 需要泛化到更大规模的图上
  \item Quantum Volume 更大的电路需要更长的 rollout
\end{itemize}

\subsubsection{观测增强}

当前 SABRE 特征提供了 front\_layer 距离信息，但缺失了 gate 权重的区分（如 CX vs CZ 的不同噪声特征）。可以考虑加入：
\begin{itemize}
  \item 每个 front\_layer 门的执行代价（保真度影响）
  \item 门的关键路径深度（影响并行度）
  \item 量子比特的剩余活动度（未来依赖数量）
\end{itemize}

\subsubsection{硬件需求}

训练 Phase 2（100K 步，Aer 密度矩阵仿真）在 5 量子比特上需要约 5--7 GPU 小时。扩展到 10+ 量子比特时，Aer 仿真时间随量子比特数指数增长（密度矩阵 $O(4^m)$），需要：
\begin{itemize}
  \item GPU 加速（单精度仿真在 A100 上可扩展到 10--15 量子比特）
  \item 对于更大规模，需要切换到张量网络仿真或 stabilization-based 方法
  \item 或者使用训练好的保真度预测器（GNN predictor）替代 Aer 仿真
\end{itemize}

\subsubsection{迁移学习与泛化}

当前模型在训练时见过的三种拓扑上表现良好，但对未见过拓扑的零样本迁移能力尚未验证。后续研究可以探索：
\begin{itemize}
  \item 元学习（MAML）实现少样本适应新拓扑
  \item 基于图编辑距离的拓扑相似度度量
  \item 域随机化：训练时随机生成多种拓扑提高泛化性
\end{itemize}

\subsubsection{分布式训练}

当前训练是单进程的。扩展到大电路和多拓扑需要分布式 PPO：
\begin{itemize}
  \item 多个 rollout worker 并行收集体验
  \item 中心化 learner 负责参数更新
  \item 每个 worker 持有一个独立的电路和拓扑副本
\end{itemize}
